% © 2026 Ing. David Jaroš — CC BY-NC-ND 4.0
%
% This work is licensed under a Creative Commons Attribution-NonCommercial-NoDerivatives
% 4.0 International License (CC BY-NC-ND 4.0).
%
% License History: Earlier drafts (up to v0.3) were released under CC BY 4.0.
% From v0.4 onward, all material is released under CC BY-NC-ND 4.0 to protect
% the integrity of the theoretical work during ongoing academic development.
%
% See LICENSE.md for full license text.
%
% File: research_tracks/alpha_spectral/hecke_trace_B_derivation_attempt.tex
%
% Task: derive_or_reject_modular_B_from_action
% Gap ID: G137-B
% Purpose: Test whether the modular coefficient B(p)=mu(Gamma_0(p))/3=(p+1)/3
% can be derived from the UBT action S[Theta] through a Hecke-trace or winding
% sector path-integral mechanism, or must remain a conditional ansatz.

\documentclass[12pt,a4paper]{article}
\usepackage[a4paper, margin=2.5cm]{geometry}
\usepackage{amsmath,amssymb,amsthm}
\usepackage{mathtools}
\usepackage{hyperref}
\usepackage{xcolor}
\usepackage{booktabs}
\usepackage{mdframed}

\newtheorem{proposition}{Proposition}[section]
\newtheorem{lemma}[proposition]{Lemma}
\theoremstyle{remark}
\newtheorem{remark}[proposition]{Remark}

\newmdenv[backgroundcolor=green!10,linecolor=green!50!black,linewidth=1.2pt]{resultbox}
\newmdenv[backgroundcolor=yellow!12,linecolor=orange!80!black,linewidth=1.2pt]{gapbox}
\newmdenv[backgroundcolor=red!8,linecolor=red!55!black,linewidth=1.2pt]{warnbox}
\newmdenv[backgroundcolor=blue!6,linecolor=blue!55!black,linewidth=1.2pt]{insightbox}

\newcommand{\PROVED}{\textcolor{green!50!black}{\textbf{[PROVED]}}}
\newcommand{\OPEN}{\textcolor{red!70!black}{\textbf{[OPEN]}}}
\newcommand{\COND}{\textcolor{blue!70!black}{\textbf{[CONDITIONAL]}}}
\newcommand{\OBS}{\textcolor{orange!80!black}{\textbf{[OBS]}}}

\title{\textbf{Hecke-Trace Attempt for the Modular $B$-Coefficient}\\[3pt]
\large Can $B(p)=\mu(\Gamma_0(p))/3=(p+1)/3$ Be Derived from $S[\Theta]$?}
\author{Ing.\ David Jaro\v{s}}
\date{2026-05-08}

\begin{document}
\maketitle

\begin{abstract}
This note tests the strongest currently available route for deriving the modular
coefficient
\[
  B(p)=\frac{\mu(\Gamma_0(p))}{3}=\frac{p+1}{3}
\]
from the UBT action $S[\Theta]$, with focus on a Hecke-trace interpretation of
the winding-sector path integral.  We isolate three questions: (i) can the
factor $p+1$ arise as a true degeneracy from coset enumeration or Hecke traces,
(ii) can the division by $3$ be obtained from the geometry of
$\mathrm{SL}(2,\mathbb{Z})\backslash\mathbb{H}$ or from the three-dimensional
imaginary quaternion sector, and (iii) can both statements be deduced from the
action without using observed $\alpha$ or fitting to $137$.  Result: the factor
$p+1$ is an exact arithmetic invariant of $\Gamma_0(p)$ and a plausible
candidate degeneracy, but no path-integral derivation from $S[\Theta]$ is known;
the division by $3$ is explained geometrically as the normalisation
$\operatorname{vol}(\mathrm{SL}(2,\mathbb{Z})\backslash\mathbb{H})=\pi/3$, while
the alternative ``$\mathrm{Im}(\mathbb{H})$ / theta-dimension'' explanation
remains heuristic.  Therefore $B(p)=(p+1)/3$ remains a \COND\ ansatz.
\end{abstract}

\tableofcontents
\newpage

\section{Problem statement}

The winding-sector effective potential used in the T3\_ALPHA route is
\begin{equation}
  V_{\mathrm{eff}}(n)=n^2-B\,n\ln n,
  \qquad n\in \mathbb{Z}_{>0}.
  \label{eq:veff}
\end{equation}
To place the minimum at a prime $n^*=p$, one requires
\begin{equation}
  2p = B(\ln p + 1).
  \label{eq:stationary}
\end{equation}
The modular ansatz proposes
\begin{equation}
  B(p)=\frac{\mu(\Gamma_0(p))}{3}=\frac{p+1}{3},
  \label{eq:bmod}
\end{equation}
which for $p=137$ gives $B(137)=46$.

\begin{warnbox}
\textbf{Hard rule.}  No observed value of $\alpha$ is used anywhere in this
note.  The question is only whether~\eqref{eq:bmod} follows from the UBT action
$S[\Theta]$.
\end{warnbox}

\section{What is exact before physics enters}

\subsection{Coset count and modular volume}

For prime level $p$, the standard congruence-subgroup formula gives
\begin{equation}
  \mu(\Gamma_0(p)) = [\mathrm{SL}(2,\mathbb{Z}):\Gamma_0(p)] = p+1.
  \label{eq:index}
\end{equation}
Equivalently, the right cosets are in bijection with the projective line
$\mathbb{P}^1(\mathbb{F}_p)$, so
\begin{equation}
  |\Gamma_0(p)\backslash \mathrm{SL}(2,\mathbb{Z})|
  = |\mathbb{P}^1(\mathbb{F}_p)|
  = p+1.
  \label{eq:p1}
\end{equation}
The hyperbolic area of $X_0(p)$ is then
\begin{equation}
  \operatorname{vol}(X_0(p))
  = \frac{\pi}{3}\,\mu(\Gamma_0(p))
  = \frac{\pi(p+1)}{3},
  \label{eq:vol}
\end{equation}
and therefore
\begin{equation}
  \frac{\operatorname{vol}(X_0(p))}{\pi} = \frac{p+1}{3}.
  \label{eq:volnorm}
\end{equation}

\begin{resultbox}
\textbf{Exact arithmetic content.}  The factor $p+1$ is exact.  The division by
$3$ is exact.  But these are modular-curve identities, not yet consequences of
$S[\Theta]$.
\end{resultbox}

\subsection{What the Hecke-trace idea would need}

The candidate mechanism is to rewrite the one-loop determinant on the
complex-time torus in a basis adapted to Hecke correspondences, schematically
\begin{equation}
  \Gamma_{\mathrm{1-loop}}[\bar{\Theta}]
  \stackrel{?}{=}
  -\frac{1}{2}\log\det{}'(\nabla^\dagger\nabla)_{\bar{\Theta}}
  \stackrel{?}{=}
  \sum_{p} c_p \,\mathrm{Tr}(T_p) + \cdots,
  \label{eq:hecke-schematic}
\end{equation}
with $\mathrm{Tr}(T_p)$ producing a factor proportional to $p+1$ on a suitable
space of winding states.  A second possible route is that the path integral over
the level-$p$ winding sector decomposes into $p+1$ equal-weight saddles labelled
by cosets of $\Gamma_0(p)$.

\begin{insightbox}
The current evidence supports the \emph{possibility} that $p+1$ is a degeneracy
factor, but there is no explicit UBT calculation establishing the required
operator, measure, or saddle decomposition.
\end{insightbox}

\section{Attempt 1: derive $p+1$ from a Hecke trace}

\subsection{Minimal mechanism}

To obtain~\eqref{eq:bmod} from~\eqref{eq:hecke-schematic}, the following chain
would be required:
\begin{enumerate}
  \item The complex-time sector of $S[\Theta]$ must define a modularly covariant
        one-loop operator on $\tau\in\mathbb{H}$.
  \item The spectrum of that operator must decompose into Hecke orbits at prime
        level $p$.
  \item The weighted trace must reduce to a uniform orbit count,
        $\mathrm{Tr}(T_p)\propto p+1$.
  \item The overall normalisation must be exactly $1/3$.
\end{enumerate}

\subsection{Where the attempt stops}

\begin{proposition}
From the current repository derivations, one cannot deduce the existence of an
operator on the UBT winding sector whose trace is exactly $\mu(\Gamma_0(p))$.
\end{proposition}

\begin{proof}
The existing UBT derivations prove the shape $V_{\mathrm{eff}}(n)=n^2-Bn\ln n$
and prove the modular identity $\mu(\Gamma_0(p))=p+1$.  They do \emph{not}
construct:
\begin{itemize}
  \item a Hecke-equivariant Hilbert space of winding states,
  \item a UBT path-integral measure reduced to a finite set of $p+1$ equal
        saddle classes, or
  \item a determinant formula whose coefficient is the trace of $T_p$.
\end{itemize}
Hence the passage from arithmetic identity to action-level coefficient is
missing.
\end{proof}

\begin{gapbox}
\textbf{Gap H1.}  The Hecke-trace mechanism is a candidate only.  No explicit
trace formula of the form
\[
  B(p)=\frac{1}{3}\,\mathrm{Tr}(T_p)
\]
has been derived from $S[\Theta]$.
\end{gapbox}

\section{Attempt 2: derive $p+1$ from winding-sector degeneracy}

\subsection{Best possible statement}

If the winding-level-$p$ sector split into $\mu(\Gamma_0(p))=p+1$ inequivalent,
equal-action saddles, then the one-loop contribution could schematically take
the form
\begin{equation}
  Z_p
  \stackrel{?}{=}
  \sum_{a\in \Gamma_0(p)\backslash\mathrm{SL}(2,\mathbb{Z})}
  e^{-S_{\mathrm{eff}}[\Theta_a]}
  = (p+1)e^{-S_*}
  \label{eq:degeneracy}
\end{equation}
and therefore yield a multiplicity factor $p+1$ in the effective coupling.

\subsection{Why this is not yet a derivation}

Equation~\eqref{eq:degeneracy} requires two non-trivial ingredients that are not
presently proved:
\begin{enumerate}
  \item \textbf{Enumeration}: the relevant UBT winding vacua must truly be
        indexed by $\Gamma_0(p)\backslash\mathrm{SL}(2,\mathbb{Z})$.
  \item \textbf{Uniform weight}: all such vacua must contribute with the same
        one-loop action.
\end{enumerate}
Neither statement follows from the current action-level analysis.  The existing
one-loop and KK+winding calculations produce the correct $n\ln n$ structure but
not the coset multiplicity.

\begin{gapbox}
\textbf{Gap H2.}  The factor $p+1$ can be \emph{interpreted} as a possible
degeneracy, but no winding-sector path integral has been evaluated that proves
this interpretation.
\end{gapbox}

\section{Attempt 3: explain the division by $3$}

\subsection{Robust explanation: modular-area normalisation}

The exact source of the denominator in~\eqref{eq:bmod} is the fundamental-domain
volume
\begin{equation}
  \operatorname{vol}(\mathrm{SL}(2,\mathbb{Z})\backslash\mathbb{H})=\frac{\pi}{3}.
  \label{eq:fundvol}
\end{equation}
Thus
\begin{equation}
  \frac{\mu(\Gamma_0(p))}{3}
  =
  \frac{\operatorname{vol}(X_0(p))}{\pi}.
  \label{eq:three-volume}
\end{equation}
This is a clean geometric normalisation and does not require any extra physical
assumption.

\subsection{Heuristic explanation: $\mathrm{Im}(\mathbb{H})$ / theta dimension}

Several internal notes suggest a structural analogy in which the denominator $3$
is associated with the three-dimensional imaginary quaternion sector
$\mathrm{Im}(\mathbb{H})\cong\mathbb{R}^3$, or with $\vartheta_3(\tau)^3$ mode
counting.  This may be suggestive, but it does not currently follow from an
action computation.

\begin{proposition}
The division by $3$ is presently derived only as a modular-volume normalisation,
not from the UBT theta-dimension or directly from $\mathrm{Im}(\mathbb{H})$.
\end{proposition}

\begin{proof}
The modular identity~\eqref{eq:three-volume} is exact.  By contrast, the
identification
\[
  \frac{\mu(\Gamma_0(p))}{3}
  \stackrel{?}{=}
  \frac{\operatorname{vol}(X_0(p))/\pi}{\dim_{\mathbb{R}}\mathrm{Im}(\mathbb{H})}
\]
is stated in current notes only as a structural analogy; no one-loop determinant
of $\nabla^\dagger\nabla$ has been computed that outputs this factor from the
quaternionic sector itself.
\end{proof}

\begin{gapbox}
\textbf{Gap H3.}  The denominator $3$ is mathematically understood, but its UBT
action-level origin remains open.  The $\mathrm{Im}(\mathbb{H})$ explanation is
not yet stronger than a heuristic mnemonic.
\end{gapbox}

\section{Synthesis}

\begin{center}
\begin{tabular}{@{}p{0.33\linewidth}p{0.22\linewidth}p{0.33\linewidth}@{}}
\toprule
Question & Status & Reason \\
\midrule
Does coset enumeration give $p+1$ exactly? & \PROVED\ arithmetically & Yes, as
$[\mathrm{SL}(2,\mathbb{Z}):\Gamma_0(p)]=p+1$ \\
\addlinespace
Does $p+1$ appear as a UBT path-integral degeneracy? & \OPEN & No explicit
winding-sector path integral or saddle count exists \\
\addlinespace
Can a Hecke trace produce $p+1$? & \OPEN & Plausible candidate, but no
Hecke-equivariant UBT operator has been derived \\
\addlinespace
Is the division by $3$ exact? & \PROVED\ geometrically & It is the modular-area
normalisation $\operatorname{vol}(\mathrm{SL}(2,\mathbb{Z})\backslash\mathbb{H})=\pi/3$ \\
\addlinespace
Is the division by $3$ derived from $\mathrm{Im}(\mathbb{H})$ or theta dimension? & \OPEN & Existing argument is heuristic only \\
\addlinespace
Can $B(p)=(p+1)/3$ be claimed as derived from $S[\Theta]$? & \textbf{\OPEN} &
The bridge from arithmetic invariant to action coefficient is missing \\
\bottomrule
\end{tabular}
\end{center}

\section{Final verdict}

\begin{resultbox}
\textbf{Verdict.}  The modular coefficient
\[
  B(p)=\frac{\mu(\Gamma_0(p))}{3}=\frac{p+1}{3}
\]
\emph{cannot currently be derived from the UBT action $S[\Theta]$}.  The factor
$p+1$ is an exact arithmetic invariant and a viable candidate degeneracy from
either Hecke traces or winding-sector cosets, but neither mechanism has been
realised in a first-principles path-integral computation.  The denominator $3$
is exactly explained by modular-area normalisation, while the
``$\mathrm{Im}(\mathbb{H})$ / theta-dimension'' explanation remains heuristic.
Therefore $B(p)$ must remain a \COND\ modular ansatz.
\end{resultbox}

\section{Implication for Gap G137-B}

Gap G137-B is \emph{not closed} by the present analysis.  The honest statement
for the T3\_ALPHA route is:
\begin{quote}
The formula $B(p)=(p+1)/3$ is structurally motivated by modular geometry and is
compatible with the prime-attractor framework, but it is not derived from
$S[\Theta]$ and therefore remains conditional.
\end{quote}

\section*{Internal references}

\begin{itemize}
  \item \texttt{research\_tracks/alpha\_spectral/b\_coefficient\_gap\_resolution.tex}
  \item \texttt{reports/B\_gap\_final\_verdict.md}
  \item \texttt{reports/gamma0\_137\_invariants.md}
  \item \texttt{canonical/alpha/modular\_prime\_attractor\_theorem.tex}
  \item \texttt{reports/f137\_projective\_geometry\_check.md}
\end{itemize}

\end{document}
